
%%==================================================================
\section{Introduction}\label{sec:intro}
%%==================================================================

\subsection{The problem}
The class $\Pc$ consists of the languages decided by a deterministic Turing
machine in time polynomial in the input length, and $\NP$ of those decided by a
nondeterministic one in polynomial time, equivalently those whose members have
certificates of polynomial length that can be checked in polynomial time. The
question whether $\Pc=\NP$ was isolated by Cook \cite{Coo71} and Levin
\cite{Lev73}, who showed that satisfiability of propositional formulas is
$\NP$-complete; Karp \cite{Kar72} then showed that twenty-one combinatorial
problems share this property. The question had been anticipated by Edmonds
\cite{Edm65}, who proposed polynomial running time as the mark of an efficient
algorithm, and by G\"odel in a letter of 1956 to von Neumann, discussed by Sipser
\cite{Sip92}. It is one of the seven Millennium Prize Problems of the Clay
Mathematics Institute \cite{Clay}, and the surveys of Sipser \cite{Sip92},
Fortnow \cite{For09}, Aaronson \cite{Aar16} and Wigderson \cite{Wig19} describe
its place in mathematics.\footnote{Gasarch's third poll of researchers
\cite{Gas19}, taken in 2019, found that 88\% of those who answered expected
$\Pc\ne\NP$, against 83\% in 2012 and 61\% in 2002.}

No proof that $\Pc\ne\NP$, or that $\Pc=\NP$, is known, and this paper does not
give one. Its subject is the strongest kind of statement that is currently
within reach: lower bounds for explicit languages against restricted models of
computation. The most direct route to $\Pc\ne\NP$ is to show that some language
in $\NP$ has no Boolean circuits of polynomial size. Shannon \cite{Sha49}
showed that almost every Boolean function of $n$ inputs needs at least
$2^n/(2n)$ gates of fan-in two, yet for languages in $\Pc$, and even in $\NP$,
the best lower bound known is $3.1n-o(n)$ \cite{LY22}, linear in the input
length $n$.\footnote{Throughout, an \emph{explicit} function is one computed
in polynomial time, or at least in exponential time with an $\NP$ oracle; a
lower bound for a function defined by counting, as in Shannon's argument, does
not name the function.} Climbing to $\ENP$, the class of languages decided in
time $2^{O(n)}$ with an oracle for $\NP$, does not help by itself: before the
preprint \cite{OAI26} discussed below, no superlinear lower bound was known for
$\ENP$ even for circuits of fan-in-two gates, and against circuits of larger
gates and unbounded depth nothing superlinear in the number of wires was known
for any language in $\ENP$.

\subsection{What is proved}
In October 2026 \oai{} released a preprint \cite{OAI26} with a deterministic
algorithm that counts the satisfying assignments of a $k$-CNF with $n$ variables
in time $2^{O(n\log\log n/\sqrt{\log n})}$, refuting the exponential time
hypothesis of Impagliazzo and Paturi \cite{IP01}. The algorithm rests on a
contraction theorem for Boolean tensor networks \cite[Theorem D.1]{OAI26},
stated here as Theorem~\ref{thm:contraction}, and \cite[Corollary 1.4]{OAI26}
combines it with Williams' argument \cite{Wil13} to show that a language in
$\ENP$ has no circuits of $c\,n\sqrt{\log n}$ fan-in-two gates. The results
below extend both from fan-in two to gates of large fan-in. Write $\wires(C)$
for the number of wires of a circuit $C$, the sum of the fan-ins of its gates,
and $s(C)=m+\wires(C)$ when $C$ has $m$ inputs.

For integers $F,\beta\ge 1$ let $\cK(F,\beta)$ be the circuits whose gates are of
two kinds (Definition~\ref{def:gates}):
\begin{itemize}[leftmargin=2em]
\item \emph{table gates} computing an arbitrary Boolean function of at most $F$
inputs, given by its truth table;
\item \emph{weighted gates} of any fan-in $f$, computing $\varphi(\sum_i
w_iy_i)$ for integer weights $|w_i|\le 2^{\beta}$, where $\varphi$ is a
threshold $[v\ge\theta]$, an exact test $[v=\theta]$, a congruence
$[v\equiv\theta\bmod\mu]$ with $1\le\mu\le2^\beta$, or, when every weight is
one, an arbitrary function of $v\in\{0,\dots,f\}$ given by its table.
\end{itemize}
Gates of fan-in two, unbounded fan-in AND, OR and parity gates, majority and
threshold gates with weights below $2^\beta$, and $\mathrm{MOD}_\mu$ gates are
all of these kinds, in any combination and at any depth.

\begin{theoremA}[Counting, Theorem~\ref{thm:counting}]
There is a deterministic algorithm that, given a circuit $C\in\cK(s^{4/5},s^{2/3})$
with $s=s(C)\ge 2$, described by $L$ bits, computes the number of inputs on which
$C$ outputs one in time $(s+L)^{O(1)}2^{O(s/\sqrt{\log s})}$ and space
$(s+L)^{O(1)}$. Since $L\le s^{O(1)}2^{s^{4/5}}$, the time is
$2^{O(s/\sqrt{\log s})}$.
\end{theoremA}

\begin{theoremB}[Lower bound, Theorem~\ref{thm:lower}]
There are a language $\Wl\in\ENP$ and a constant $c>0$ such that, for
infinitely many $N$, every circuit in $\cK(N^{3/4},N^{1/2})$ with $N$ inputs
that decides $\Wl$ on inputs of length $N$ has more than $cN\sqrt{\log N}$ wires.
\end{theoremB}

A fan-in-two circuit with $g$ gates has at most $2g$ wires, so Theorem~B
contains \cite[Corollary 1.4]{OAI26}. Table gates of fan-in $N^{3/4}$ cannot be
replaced by fan-in-two gates at linear cost: most functions of $k=N^{3/4}$
inputs need at least $2^k/(2k)$ such gates \cite{Sha49}. A threshold gate of
fan-in $f$ with weights of $\beta$ bits is computed by the usual adder circuits
with $O(f\beta)$ fan-in-two gates, which for $\beta=N^{1/2}$ is far from linear
in its $f$ wires. Theorem~A therefore works with the gates themselves rather
than through a simulation.

The counting theorem also applies to integer programs. For a matrix $A$ write
$\nu(A)$ for its number of nonzero entries.

\begin{theoremC}[Zero-one programs, Corollary~\ref{cor:ilp}]
Let $A\in\Z^{r\times n}$ and $b\in\Z^r$, and put $\sigma=n+\nu(A)$. If every
entry of $A$ has absolute value at most $2^{\sigma^{2/3}}$, the number of
$x\in\{0,1\}^n$ with $Ax\le b$ can be computed in time
$L^{O(1)}2^{O(\sigma/\sqrt{\log\sigma})}$, and a linear objective with
coefficients of the same size can be maximised over these points within the
same bound.
\end{theoremC}

For sparse programs, with $\nu(A)=O(n)$, this is $2^{O(n/\sqrt{\log n})}$.
Before \cite{OAI26}, the best algorithm for zero-one programs with $O(n)$
constraints ran in time $2^{(1-\epsilon)n}$ for a constant $\epsilon>0$
depending on the ratio of constraints to variables \cite{ILPS14}; a
subexponential algorithm for sparse zero-one programs would already have refuted
the exponential time hypothesis, since $3$-SAT with $O(n)$ clauses is such a
program.

The argument that turns Theorem~A into Theorem~B works for any exponent. Its
general form is Theorem~\ref{thm:transfer}:

\begin{theoremD}[Transfer]
Let $\gamma>0$ be rational. If satisfiability of circuits in $\cK(s^{4/5},s^{2/3})$ with
$s$ wires and inputs can be decided in time $L^{O(1)}2^{O(s/(\log s)^\gamma)}$,
then some language in $\ENP$ needs more than $cN(\log N)^\gamma$ wires in
$\cK(N^{3/4},N^{1/2})$ for infinitely many $N$.
\end{theoremD}

Theorem~A supplies $\gamma=\tfrac12$, and a $2^{O(s/\log s)}$ algorithm would
give $cN\log N$. Superpolynomial lower bounds are a different matter. In
Williams' theorem \cite{Wil13} what has to be beaten is exhaustive search over
the $n$ inputs: deciding satisfiability of circuits of polynomial size with $n$
inputs in time $2^n/n^{\omega(1)}$ would show $\NEXP\not\subseteq\Ppoly$. The
bound of Theorem~A is measured in the size $s$, and it is below $2^n$ only while
$s=O(n\sqrt{\log n})$. Section~\ref{sec:limits} explains why the method stops
at $\sqrt{\log N}$ and what each further step would require; none of these
steps would separate $\Pc$ from $\NP$.

\subsection{Plan}
Part~I surveys the problem. Section~\ref{sec:history} gives definitions, history
and equivalent forms; Section~\ref{sec:barriers} the three barriers that every
proof must pass; Section~\ref{sec:circuits} the circuit lower bounds known
today; Section~\ref{sec:algorithms} the algorithmic method, the exponential time
hypothesis and its refutation; Section~\ref{sec:otherroutes} the other
programmes; and Section~\ref{sec:repositories} the repositories and the
catalogue of 6,480 arXiv papers assembled for this survey. Part~II contains the
new results: circuits and tensor networks in Section~\ref{sec:networks},
Theorem~A in Section~\ref{sec:counting}, zero-one programs in
Section~\ref{sec:ilp}, Theorems~B and~D in Section~\ref{sec:lower}, and the
limits of the method in Section~\ref{sec:limits}.
Appendix~\ref{app:contraction} reproduces the proof of the contraction theorem
of \oai{} \cite[Appendix D.4]{OAI26} on which Theorem~A rests, so that the paper
is complete, and Appendix~\ref{app:checks} describes the programs that check its
finite identities.

\subsection*{Conventions}
Logarithms are to base two unless written $\ln$. $[n]=\{1,\dots,n\}$. For a
proposition $P$, $[P]$ is $1$ if $P$ holds and $0$ otherwise. For an integer or
rational number $x\ne0$, $\vtwo(x)$ is its two-adic valuation. Running times
count bit operations on a random-access machine unless a Turing machine is named.
